If two labelled invariants are equivalent, then by definition they agree, after label-preserving bijections, in \(\mathfrak C(P)\), \(N\), \(\omega\), every comparison label, every primitive residual label, and every signed incidence.  Because \(\mathfrak C(P)=\Omega(P)\), the blockwise formula
\[
 b_k^{\mathrm{sep}}(P)
 =\omega_k\frac{\mathfrak C_{kk}(P)^{-1}\mathbf 1}
 {\mathbf 1^{\prime}\mathfrak C_{kk}(P)^{-1}\mathbf 1}
\]
is computable from the invariant.  Therefore
\[
 D([\mathfrak H(P)])
 :=N^{\prime}\mathfrak C(P)b^{\mathrm{sep}}(P)
\]
is evaluable and constant on each label-preserving equivalence class.  Its zero and nonzero inverse images are exactly \(\mathscr E_0\) and \(\mathscr E_+\).

For an invariant induced by a compatible \(Q\in\mathcal M\), Theorem \(\ref{thm:separate-joint-frontier}\) gives
\[
 V_{\mathrm{sep}}-V_{\mathrm{joint}}
 =D([\mathfrak H(P)])^{\prime}
  (N^{\prime}\Omega N)^{-1}
  D([\mathfrak H(P)]).
\]
The middle matrix is positive definite whenever the contrast space is nonzero; in the zero-dimensional case \(D=0\).  Thus the gap is zero on \(\mathscr E_0\) and strictly positive on \(\mathscr E_+\).  This is exactly an evaluable restatement of the frontier, not an independent structural classification.

For nonemptiness, apply Theorem \(\ref{thm:open-strict-gap-family}\) with
\[
 \rho=\frac12,
 \qquad
 \delta=\frac1{16}.
\]
These values satisfy \(0<|\rho|<\sqrt7/3\) and \(0<|\delta|<1/8\).  That theorem supplies a bounded compatible law in \(\mathcal M\) with the stated architecture and
\[
 \Omega=
 \begin{pmatrix}
 3/8&3/16&0\\
 3/16&3/8&3/32\\
 0&3/32&3/8
 \end{pmatrix},
 \qquad
 b^{\mathrm{sep}}=
 \begin{pmatrix}\omega_1/2\\ \omega_1/2\\ \omega_2\end{pmatrix}.
\]
Its exact gap is
\[
 \frac{3(1/2)^2\omega_2^2}{32}
 =\frac{3\omega_2^2}{128}>0.
\]
Hence \(D\ne0\), the induced invariant belongs to \(\mathscr E_+\), and the strict class is nonempty.  The same theorem proves that this witness is a point of an open family indexed by \((\rho,\delta)\), with nonzero effect \(\delta\).